% Clase 11 --- Los tres campos dentro del dieléctrico (§1.3).
% Clase 10 ya mostró las moléculas orientándose; acá lo que importa NO son las
% moléculas sino el campo que producen sus cargas de superficie: E' es colineal
% con E0 (por isotropía no hay otra dirección disponible) y opuesto, así que el
% campo total queda MENOR que el aplicado. Eso es todo lo que hace falta para
% que K_E >= 1.
% Se respeta la orientación del pizarrón de esta clase: -Q arriba, +Q abajo,
% de modo que E0 apunta hacia arriba.
\begin{center}
\begin{tikzpicture}[scale=1]

  % ---- placas del sistema ----
  \draw[figblue, line width=1.6pt] (-2.3,2.1) -- (2.3,2.1);
  \foreach \x in {-2.0,-1.4,...,2.1} {\node[figlbl,figblue,scale=0.7] at (\x,2.3) {$-$};}
  \node[figlbl, figblue, anchor=west] at (2.45,2.25) {$-Q$};

  \draw[figred, line width=1.6pt] (-2.3,-2.1) -- (2.3,-2.1);
  \foreach \x in {-2.0,-1.4,...,2.1} {\node[figlbl,figred,scale=0.7] at (\x,-2.3) {$+$};}
  \node[figlbl, figred, anchor=west] at (2.45,-2.25) {$+Q$};

  % ---- dieléctrico ----
  \fill[figgray!8] (-2.15,-1.65) rectangle (2.15,1.65);
  \draw[figgray, line width=0.7pt] (-2.15,-1.65) rectangle (2.15,1.65);
  \fill[figred!14] (-2.15,1.35) rectangle (2.15,1.65);
  \foreach \x in {-1.8,-1.2,...,1.9} {\node[figlbl,figred,scale=0.7] at (\x,1.5) {$+$};}
  \fill[figblue!14] (-2.15,-1.65) rectangle (2.15,-1.35);
  \foreach \x in {-1.8,-1.2,...,1.9} {\node[figlbl,figblue,scale=0.7] at (\x,-1.5) {$-$};}
  \node[figlblaux, anchor=west, align=left] at (2.45,1.5)
    {cargas de\\[-0.15em] polarización};
  \node[figlblaux, anchor=west] at (2.45,-1.5) {$K_E$};

  % ---- los tres campos, en columnas separadas y rotulados al costado ----
  \draw[figvecres] (-1.7,-1.15) -- (-1.7,1.15);
  \node[figlbl, figred, anchor=west] at (-1.6,0.75) {$\vec E_0$};

  \draw[figvecaux, figblue, line width=1pt,
        -{Latex[length=2.2mm,width=1.6mm]}] (-0.25,0.5) -- (-0.25,-0.5);
  \node[figlbl, figblue, anchor=west] at (-0.15,-0.28) {$\vec E\,'$};

  \draw[figvecres, line width=1.3pt] (1.25,-0.7) -- (1.25,0.7);
  \node[figlbl, figred, anchor=west] at (1.35,0.4) {$\vec E$};

  % ---- suma vectorial, al costado ----
  \begin{scope}[xshift=7.9cm, yshift=-0.1cm]
    \draw[figvecres] (0,-1.15) -- (0,1.15);
    \node[figlbl, figred, anchor=south, align=center] at (0,1.2)
      {$E_0$\\[-0.25em] \footnotesize cargas libres};
    \draw[figvecaux, figblue, line width=1pt,
          -{Latex[length=2.2mm,width=1.6mm]}] (2.1,0.45) -- (2.1,-0.45);
    \node[figlbl, figblue, anchor=south, align=center] at (2.1,1.2)
      {$E'$\\[-0.25em] \footnotesize polarización};
    \draw[figvecres, line width=1.3pt] (4.2,-0.7) -- (4.2,0.7);
    \node[figlbl, figred, anchor=south, align=center] at (4.2,1.2)
      {$E$\\[-0.25em] \footnotesize total};
    \node[figlbl, anchor=north, align=center] at (2.1,-1.55)
      {$E = E_0 - E' < E_0$
       \ $\Longrightarrow$ \ $K_E \geq 1$};
  \end{scope}

\end{tikzpicture}\\[0.5em]
{\small\itshape Las dos propiedades de $\vec E\,'$ no se postulan, se deducen de
las hipótesis sobre el material. Que sea \textbf{colineal} con $\vec E_0$ es
consecuencia de la \textbf{isotropía}: la única dirección privilegiada del
problema es la que impone el campo aplicado, así que no hay ninguna otra hacia
la cual $\vec E\,'$ pudiera apuntar. Que sea \textbf{opuesto} es la disposición
de las cargas de polarización, que quedan de signo contrario a la placa vecina.
Como los dos vectores se restan, el campo dentro del material es siempre
\emph{menor} que el aplicado ---y eso, y no otra cosa, es lo que obliga a que
$K_E \geq 1$---. El dieléctrico no cancela el campo como haría un conductor:
sólo lo debilita en un factor $K_E$.}
\end{center}
